\chapter{Машина з живих нейронів}
\chaptermark{Машина з живих нейронів}
\label{sec:livingmachine}

\section{Стенд із відкритою схемою}

У розділі~\ref{sec:debugging} ми налагоджували машину, в якої немає схеми: щуп чує тисячу нейронів із вісімдесяти мільярдів, а все інше доводиться вгадувати. Інженер у такому становищі зробив би просту річ --- зібрав би невеликий шматок схеми на макетній платі, де видно кожну деталь. З нейронами це теж можна зробити.

Нейрони кори розділяють і вирощують на чипі з ґраткою електродів --- від кількох десятків до десятків тисяч. За кілька тижнів клітини проростають відростками й з'єднуються в мережу з тисяч або сотень тисяч нейронів, здебільшого \nt{GLUT} і меншої частки \nt{GABA}, яка залежить від препарату: у культурах гіпокампа це близько $6\%$ нейронів~\cite{Benson1994}. Кожен електрод уміє і слухати, як щуп розділу~\ref{sec:debugging}, і подавати струм, як тестовий сигнал. Другий різновид стенда --- \emph{органоїди}: тривимірні грудочки нервової тканини розміром близько міліметра, вирощені зі стовбурових клітин людини~\cite{Lancaster2013}. На таких стендах живі мережі працюють місяцями; є майданчики, де досліди на органоїдах можна ставити віддалено, через мережу, понад сто днів поспіль~\cite{Jordan2024}. Цю галузь називають \emph{біологічними обчисленнями} або \emph{органоїдним інтелектом}~\cite{Smirnova2023}.

Стенд має дві важливі відмінності від мозку. Він крихітний: нейронів у ньому в сто тисяч разів менше. І в ньому немає широкомовних каналів: нейронів \nt{DOPA}, \nt{SERO}, \nt{NORA} і \nt{ACET} у культурі з кори немає. Це машина, яка має шину даних і керування потоком, але не має регістрів керування. Подивімося, на що вона здатна.

\section{Що мережа робить сама}

Якщо культуру залишити в спокої, вона не мовчить і не шумить рівномірно. Час від часу майже вся мережа спалахує разом --- \emph{залпом} --- і знову затихає; інтервали між залпами --- від секунди до хвилин, залежно від культури~\cite{Wagenaar2006}. Для інженера це знайома картина: підсилювач із сильним додатним зворотним зв'язком і без навантаження перетворюється на генератор.

Механізм ми вже знаємо. Сильні взаємні зв'язки \nt{GLUT}-нейронів запускають лавину, як у розділі~\ref{sec:gaba} при порушенні умов твердження~\ref{prop:ei}. Лавина швидко виснажує запас медіатора в синапсах~\eqref{eq:depression}, і мережа замовкає. Запас відновлюється зі сталою часу $\tau_{\text{відн}}$, і коли відновиться частка $x^*$, достатня для нової лавини, настає новий залп. Інтервал між залпами
\begin{empheq}[box=\keybox]{equation}
T_{\text{залп}} \;\approx\; \tau_{\text{відн}}\,\ln\frac{1}{1-x^*} .
\label{eq:burstinterval}
\end{empheq}
При $\tau_{\text{відн}}=0.8\,\text{с}$ із розділу~\ref{sec:code} і $x^*=0.9$ виходить близько двох секунд, при $x^*=0.99$ --- близько чотирьох. Це оцінка моделі з $\tau_{\text{відн}}$ книги, і вона потрапляє в короткий край виміряних інтервалів. Пряме вимірювання на мікрокультурах показує, що запас медіатора після залпу відновлюється за десятки секунд~\cite{CohenSegal2011}, тобто $\tau_{\text{відн}}$ там більша; з такою $\tau_{\text{відн}}$ формула~\eqref{eq:burstinterval} дає десятки секунд і хвилини, як у повільних культурах. Це генератор релаксаційного типу, родич генератора ритму з розділу~\ref{sec:muscles}: там його перезаряджала втома груп, тут --- запас медіатора.

Залпи --- це те, що мережа робить, коли їй нічого робити. У живому мозку схожа картина буває в глибокому сні (розділ~\ref{sec:sleep}) і в мозку, що розвивається, доки до нього не надійшли справжні входи. Ця схожість наводить на думку, але однаковість механізмів --- гіпотеза.

\section{Як навчати мережу без учителя-дофаміну}

Оскільки \nt{DOPA} в культурі немає, сигнал «краще чи гірше» доводиться подавати самим --- електродами. Досліди показують, що мережа на стенді справді змінює відповіді, і найцікавіше в них --- \emph{чим} її навчають.

\emph{Учитель, який замовкає.} Мережу стимулюють через один електрод раз на одну--три секунди, доки на іншому електроді не з'явиться потрібна відповідь через $50\pm10\ms$ після стимулу. Щойно вона з'явилася, стимуляцію вимикають на п'ять хвилин. Після кількох таких циклів потрібна відповідь з'являється на стимул одразу~\cite{Shahaf2001}. Винагорода тут --- тиша: стимул розхитує мережу, а припинення стимулу залишає її в тому стані, який дав правильну відповідь. Це спуск за градієнтом через випадкові відхилення з розділу~\ref{sec:dofa} (твердження~\ref{prop:gradient}), де роль помилки відіграє сам факт трусіння: трусять, доки не стане правильно.

\emph{Мережа, яка грає в теніс.} У досліді на стенді зі щільною ґраткою електродів близько мільйона нейронів підключили до комп'ютерної гри в теніс з однією ракеткою~\cite{Kagan2022}. Положення м'яча подавали струмом на «чутливі» електроди: де м'яч --- кодом місця, як далеко --- частотою. Активність у двох «рухових» ділянках рухала ракетку вгору або вниз. Правило навчання було незвичним. Якщо ракетка відбила м'яч, мережа отримувала короткий \emph{передбачуваний} стимул; якщо схибила --- кілька секунд \emph{непередбачуваного} випадкового струму. За двадцять хвилин гри серії відбитих м'ячів подовжувалися. Значуще зростання всередині сеансу було тільки за повного зворотного зв'язку; коли замість стимулу була тиша, культура наприкінці сеансу теж грала краще, ніж без зворотного зв'язку, але з меншим ефектом, а стійкого навчання від сеансу до сеансу впродовж кількох днів не було. Докладний розбір цього досліду --- у розділі~\ref{sec:dishbrain}.

Автори пояснили результат гіпотезою, що мережа діє так, щоб її вхід був передбачуваним~\cite{Friston2010}. Це та сама ідея, що й економія імпульсів у розділі~\ref{sec:code} та передбачення кадру в розділі~\ref{sec:recognition}: машина витрачає менше, коли вхід очікуваний. Але і сам дослід, і особливо слова авторів про «розумність» культури викликали різку критику: ефект невеликий, а пояснити його можна й простіше~\cite{Balci2023}. Чесна оцінка така: культура на стенді змінює поведінку під дією зворотного зв'язку --- це виміряно, а те, що вона вчиться саме передбачати свій вхід, --- гіпотеза.

\emph{Мережа, яка керує тілом.} Подібно до цього культуру підключали до простого віртуального «тіла» і вчили його рухатися до цілі, добираючи просторово-часовий малюнок тренувальних стимулів~\cite{Bakkum2008}. У жодному з цих дослідів немає дофаміну: учителем є малюнок струму, який обирає інженер. Що змінюється, коли учителем стає програма або сам дофамін, --- у підрозділах~\ref{sec:dishdopamine}--\ref{sec:cartpole}.

\begin{figure}[t]
\centering
\begin{tikzpicture}[>=Stealth,
  blk/.style={draw,very thick,rounded corners=2pt,align=center,font=\footnotesize,minimum height=0.8cm,fill=blue!5}]
% (a) closed loop
\node[font=\small] at (1.5,4.3) {(а) замкнена петля};
\node[blk,text width=2.6cm] (G) at (1.5,3.3) {гра: м'яч і ракетка};
\draw[very thick,fill=orange!10,rounded corners=3pt] (0.5,0.2) rectangle (2.5,1.6);
\foreach \x in {0.7,1.0,...,2.35}{\foreach \y in {0.4,0.7,1.0,1.3}{\fill[gray!60] (\x,\y) circle (0.04);}}
\draw[->,thick,blue!60!black] (G.west) -| (-0.5,0.9) -- (0.5,0.9);
\node[font=\scriptsize,blue!60!black,anchor=east,align=right] at (-0.6,2.1) {м'яч\\$\to$ струм};
\draw[->,thick,red!70!black] (2.5,0.9) -- (3.5,0.9) |- (G.east);
\node[font=\scriptsize,red!70!black,anchor=west] at (3.6,2.1) {імпульси $\to$ ракетка};
\node[font=\scriptsize] at (1.5,-0.2) {нейрони на ґратці електродів};
\node[font=\scriptsize,align=center] at (1.5,-0.85) {відбив: передбачуваний стимул\\промах: випадковий струм};
% (b) reservoir
\begin{scope}[xshift=7.9cm]
\node[font=\small] at (1.6,4.3) {(б) резервуар};
\node[blk,text width=1.1cm] (X) at (-0.4,1.0) {вхід\\$u(t)$};
\draw[very thick,fill=orange!10] (1.6,1.0) circle (0.8);
\foreach \a/\r in {20/0.35,80/0.5,150/0.3,210/0.55,280/0.4,330/0.2,110/0.15}{\fill[red!60!black] ({1.6+\r*cos(\a)},{1.0+\r*sin(\a)}) circle (0.05);}
\node[font=\scriptsize] at (1.6,-0.2) {органоїд: без учителя};
\node[blk,text width=1.8cm,fill=green!8] (W) at (4.0,1.0) {зчитування\\$W_{\text{вих}}$};
\draw[->,thick] (X) -- (0.8,1.0);
\draw[->,thick] (2.4,1.0) -- node[above,font=\scriptsize] {$\mathbf r(t)$} (W);
\draw[->,thick] (W) -- (4.0,2.3) node[above,font=\scriptsize] {$y(t)$};
\node[font=\scriptsize,green!40!black] at (4.0,0.3) {навчається з учителем};
\end{scope}
\end{tikzpicture}
\caption{Два способи змусити живу мережу обчислювати. (а) Замкнена петля: положення м'яча подається струмом, імпульси рухових ділянок рухають ракетку, а передбачуваний або випадковий стимул після удару слугує учителем. (б) Резервуар~\eqref{eq:reservoir}: органоїду не подають навчального сигналу, він розщеплює вхід на безліч нелінійних відгуків із пам'яттю; за помилкою навчається лише одне лінійне зчитування ззовні. Зв'язки самого органоїду при цьому теж змінюються --- без учителя.}
\label{fig:livingmachine}
\end{figure}

\section{Живий резервуар}

Є й інший спосіб використати живу мережу: не навчати її взагалі. У техніці він називається \emph{обчисленням на резервуарі}~\cite{Maass2002,JaegerHaas2004}. Беруть готову нелінійну систему з пам'яттю --- будь-яку, аби її відгук на вхід був багатим і залежав від недавнього минулого. Вхід $u(t)$ розгойдує її, виходить вектор відгуків $\mathbf r(t)$, і навчається тільки лінійне зчитування
\begin{empheq}[box=\keybox]{equation}
y(t) \;=\; W_{\text{вих}}\,\mathbf r(t) ,
\label{eq:reservoir}
\end{empheq}
ваги якого добирають правилом найменших квадратів~\eqref{eq:lms} з розділу~\ref{sec:motorlearning}. Резервуар робить те саме, що робили маленькі \nt{GLUT}-нейрони розділу~\ref{sec:motorlearning}: розщеплює вхід на довгий вектор, у якому потрібні класи стають роздільними площиною (розділ~\ref{sec:dendrites}).

Органоїд на ґратці електродів підійшов на роль такого резервуара. Звуки мовлення, подані на нього малюнками струму, після навчання одного лінійного зчитування розрізнялися за мовцем із точністю близько $78\%$~\cite{Cai2023}. Точність $78\%$ --- цифра, яку повідомляють самі автори і яку повторює преса. Результат корисний, але важливо розуміти, де в ньому «розум»: органоїд дає нелінійність і згасну пам'ять, навчання за помилкою зроблено ззовні, у кремнії (рис.~\ref{fig:livingmachine}). При цьому органоїд не лишається незмінним: за даними авторів, під час тренування змінювалася його власна функціональна зв'язність, і це називають навчанням без учителя в самому органоїді~\cite{Cai2023}. Яку частку точності дає ця зміна, а яку --- зчитування, не розділено.

\section{Що стенд каже про машину}

\emph{Енергія.} За оцінкою~\eqref{eq:energy} імпульс коштує близько $10^{-10}$~Дж. Культура з мільйона нейронів за частоти один імпульс на секунду витрачає на передавання сигналів
\begin{empheq}[box=\keybox]{equation}
P \;\approx\; 10^{6}\times1\,\text{с}^{-1}\times10^{-10}\,\text{Дж} \;=\; 10^{-4}\,\text{Вт},
\label{eq:dishpower}
\end{empheq}
десяту частку мілівата. Електроніка, яка стимулює, записує й обробляє ці імпульси, витрачає на порядки більше. Як і в розділі~\ref{sec:debugging}, вузьке місце --- не обчислення, а зв'язок із ним.

\emph{Відсутні деталі.} Стенд показує, що вміє шина \nt{GLUT} із керуванням потоком \nt{GABA} без регістрів керування: самоорганізовуватися, генерувати залпи, змінювати відповіді під дією зворотного зв'язку і слугувати добрим резервуаром. Чого вона без них не вміє, --- самостійно вирішити, що вважати успіхом. Цей сигнал на стенді подає інженер, а в мозку, як ми бачили в розділах~\ref{sec:dofa}--\ref{sec:acet}, --- широкомовні канали.

\emph{Етика.} Органоїди вирощують із клітин людини, і що складнішими вони стають, то серйозніше постає питання, чи може така тканина щось відчувати. Інженерний погляд підказує часткову відповідь: біль у розділі~\ref{sec:pain} --- це не сигнал одного датчика, а ціла система тривоги з воротами, регулятором гучності та широкомовними каналами, яких на стенді немає. Але це міркування, а не доведення, і сама галузь пропонує вести етичну оцінку разом із дослідами, від самого початку~\cite{Smirnova2023}.

\begin{keyblock}
\begin{proposition}[Жива мережа на стенді]
\label{prop:livingmachine}
Мережа з \nt{GLUT}- і \nt{GABA}-нейронів на ґратці електродів сама породжує залпи~\eqref{eq:burstinterval}, змінює відповіді під дією зворотного зв'язку і може слугувати резервуаром~\eqref{eq:reservoir} для навченого ззовні зчитування. Усе це виміряно. Те, що така мережа вчиться за якимось загальним правилом, наприклад передбачати свій вхід, --- гіпотеза, а гучних слів про її «розумність» більшість фахівців не приймає.
\end{proposition}
\end{keyblock}

\section{Дофамін за спалахом світла}
\label{sec:dishdopamine}

Єдиний відомий нам прямий дослід, де культурі дали дофамін як учителя, поставлено на платформі з органоїдів, до якої дослідники підключаються віддалено~\cite{Jordan2024}. У рідину, що омиває органоїди, додали дофамін у світлочутливій «клітці»: зв'язана молекула не діє на рецептори. Концентрація дофаміну в клітці --- $1$~мМ. Коли треба нагородити мережу, на неї світять ультрафіолетом: клітка розпадається, і дофамін виходить в об'єм.

Петля влаштована так. Той самий стимул подають раз у раз; якщо він викликав більше імпульсів, ніж раніше, вмикається спалах і вивільняється дофамін. За $240$ повторів кількість викликаних імпульсів загалом зростала. Автори самі пишуть, що з одного такого досліду не можна зробити висновок, що зростання спричинив дофамін~\cite{Jordan2024}.

Для інженера це перша спроба зібрати в чашці правило трьох множників~\eqref{eq:threefactor}: активність відправника й отримувача залишає слід, а спільний сигнал вирішує, чи закріпити зміну. Але сигнал тут буває лише додатним: винагорода є або її немає, від'ємної помилки $\delta<0$ установка не видає. До того ж дофамін розпливається й відкачується повільно, як концентрація в~\eqref{eq:lowpass}, тож часті спалахи можуть просто підняти його загальний рівень. Щоб відрізнити від цього навчання, потрібні два контролі: стільки ж спалахів у випадкові моменти, не пов'язані з відповіддю мережі, і ті самі спалахи без дофаміну в клітці. І потрібна перевірка після відпочинку: чи залишилася зміна, коли дофамін пішов.

\section{Дофамін навчає кремній}
\label{sec:biohybrid}

У зворотному досліді живі клітини навчають електроніку~\cite{Keene2020}. Клітини, що виділяють дофамін, вирощують у мікроканалі над органічним нейроморфним елементом --- транзистором, провідність каналу якого відіграє роль ваги синапсу. Дофамін, що вийшов із клітин, окиснюється біля електрода елемента, і це змінює провідність. Пристрій відтворює і механізм прибирання медіатора зі щілини, тому вагу можна не тільки змінити надовго, а й повернути назад.

Схемотехнічно це інтегратор із витоком за хімічним входом: вага накопичує потік дофаміну, а «відкачування» задає, як швидко вона забувається, --- те саме RC-коло, що й~\eqref{eq:lowpass}. Головне тут --- напрям зв'язку. Щуп із розділу~\ref{sec:debugging} не бачить широкомовних регістрів, бо вони хімічні. Цей елемент читає саме хімічний регістр і перетворює його на електричну вагу. Для інтерфейсів мозок---комп'ютер це шлях до датчика, який чує не тільки шину даних, а й регістри керування.

\section{Органоїди з власними дофаміновими нейронами}
\label{sec:midbrainorg}

Зі стовбурових клітин людини вміють вирощувати органоїди, схожі на ту ділянку мозку, де розташовані \nt{DOPA}-нейрони. У них є робочі дофамінові нейрони, які виділяють дофамін~\cite{Jo2016}. Такі органоїди вирощують здебільшого для вивчення хвороб. Дослідів, де їхній дофамін слугував би учителем для іншої мережі, ми не знайшли.

Для машини з живих нейронів це відсутня деталь. Стенд із розділу~\ref{sec:livingmachine} --- шина даних і керування потоком без регістрів керування; тут джерело широкомовного сигналу вже є. Бракує критика: мережі, яка обчислювала б помилку передбачення~\eqref{eq:ctd} і керувала б цими нейронами. З'єднати органоїд кори з таким органоїдом так, щоб дофамін виділявся за помилкою, --- відкрита задача, і поки що це гіпотеза, а не дослід.

\section{Органоїд балансує жердиною без дофаміну}
\label{sec:cartpole}

Найповніший із недавніх дослідів обходиться зовсім без дофаміну~\cite{Robbins2026}. Органоїди кори миші підключили замкненою петлею до задачі «жердина на візку»: активність органоїду рухає візок, положення жердини подається назад стимуляцією. Між спробами органоїд отримує короткі високочастотні навчальні стимули. Які саме --- обирає програма навчання з підкріпленням.

Результати виміряно:
\begin{itemize}
\item у більшості органоїдів навчальні сигнали, обрані програмою, дали кращий результат, ніж випадково обрані сигнали або відсутність сигналу;
\item після $45$ хвилин відпочинку покращення не зберігалося;
\item блокада рецепторів глутамату обох типів, AMPA і NMDA, скасовувала покращення.
\end{itemize}

Останній пункт показує, де живе вивчене: у \nt{GLUT}-синапсах, і через той рецептор, який у підрозділі~\ref{sec:weightstore} працює як детектор збігу входу й виходу. Другий пункт показує, чого бракує: швидка зміна є, закріплення немає, як у швидкого учня з розділу~\ref{sec:motorlearning} без повільного. А роль учителя змінилася: інженера, який обирає малюнок струму, замінила програма, але учитель, як і раніше, стоїть поза чашкою.

Серед раніших дослідів навчання культур на електродних матрицях ми також не знайшли жодного, де використовувався б дофамін: навчальним сигналом скрізь слугувала електрична стимуляція.

\section{Хто навчає культуру}

\begin{table}[t]
\centering
\footnotesize
\setlength{\tabcolsep}{4pt}
\renewcommand{\arraystretch}{1.15}
\begin{tabular}{>{\raggedright}p{2.7cm}>{\raggedright}p{2.5cm}>{\raggedright}p{3.2cm}>{\raggedright\arraybackslash}p{4.4cm}}
\toprule
Дослід & Хто навчає & Навчальний сигнал & Що відомо \\
\midrule
зупинка стимуляції за потрібної відповіді & інженер & кінець стимуляції & відповідь закріплюється --- виміряно \\
DishBrain (розд.~\ref{sec:dishbrain}) & інженер & передбачуваний або випадковий струм & значуще зростання розіграшів за сеанс лише за повного зворотного зв'язку, за тиші ефект менший --- виміряно; від сеансу до сеансу нестійкий; причина дискусійна \\
жердина на візку & програма навчання з підкріпленням & обрані програмою високочастотні стимули & краще за випадкові, але не зберігається після відпочинку --- виміряно \\
дофамін за спалахом & правило «більше імпульсів --- винагорода» & дофамін, вивільнений світлом & зростання імпульсів --- спостереження; роль дофаміну не доведено \\
біогібридний синапс & живі клітини & дофамін від клітин & тривала зміна і повернення ваги елемента --- виміряно \\
органоїди з \nt{DOPA}-нейронами & --- & --- & джерело дофаміну є; як учителя не використано \\
\bottomrule
\end{tabular}
\caption{Хто і чим навчає живі нейрони на чипі.}
\label{tab:dishteachers}
\end{table}

В усіх дослідах, де навчається сама культура, учитель поки що ззовні (таблиця~\ref{tab:dishteachers}): інженер, програма або правило, задане інженером. Дофамін потрапив у чашку лише одного разу, і його роль ще належить довести. Зібрати в чашці справжній широкомовний канал --- із критиком, помилкою і слідом, як у розділі~\ref{sec:dofa}, --- наступний крок, якого ще ніхто не зробив.


\section{Підсумок}

\begin{table}[t]
\centering
\footnotesize
\setlength{\tabcolsep}{4pt}
\renewcommand{\arraystretch}{1.15}
\begin{tabular}{>{\raggedright}p{2.8cm}>{\raggedright}p{4.2cm}>{\raggedright\arraybackslash}p{5.8cm}}
\toprule
Дослід & Що робить жива мережа & Що відомо \\
\midrule
мережа без входу & залпи з інтервалами від секунди до хвилин~\eqref{eq:burstinterval} & виміряно; механізм через виснаження синапсів --- загальноприйнята модель \\
учитель, який замовкає & потрібна відповідь закріплюється, коли стимуляцію вимикають & виміряно \\
гра в теніс & серії відбитих м'ячів подовжуються; значуще зростання за сеанс --- тільки за повного зворотного зв'язку & зміну поведінки виміряно; навчання від сеансу до сеансу нестійке; величина ефекту і пояснення передбаченням входу дискусійні \\
віртуальне тіло & рух до цілі за дібраного малюнка стимулів & виміряно \\
резервуар & нелінійність і пам'ять для зчитування~\eqref{eq:reservoir} & виміряно; навчання за помилкою --- ззовні, у кремнії; зв'язки органоїду теж змінюються \\
енергія & близько $10^{-4}$~Вт на мільйон нейронів~\eqref{eq:dishpower} & оцінка за~\eqref{eq:energy} \\
\bottomrule
\end{tabular}
\caption{Що показали машини з живих нейронів.}
\label{tab:livingmachine}
\end{table}

Машина з живих нейронів (таблиця~\ref{tab:livingmachine}) --- це стенд, на якому видно, що вміють шина даних і керування потоком без регістрів керування. Вони вміють багато, але не вміють самі вирішити, що добре. На стенді цю роль відіграє інженер зі своїм малюнком струму; у мозку --- ті нечисленні широкомовні проводи, яким присвячено середину книги.

\section{Задачі}

\begin{exercise}[\lvlA]
\label{ex:21:1}
\emph{Інтервал між залпами.} Залп виснажує запас до нуля, далі $\tau_{\text{відн}}\,\dd x/\dd t=1-x$, і новий залп починається при $x=x^*$; $\tau_{\text{відн}}=0.8$~с. (а)~Знайдіть інтервал при $x^*=0.9$ і $0.99$. (б)~Поріг $x^*$ випадково змінюється на $\pm0.005$ біля $0.99$. У якому діапазоні лежать інтервали?
\end{exercise}

\begin{exercise}[\lvlA]
\label{ex:21:2}
\emph{Енергія стенда.} (а)~Яку потужність дає оцінка~\eqref{eq:dishpower} для всього мозку ($8.6\cdot10^{10}$ нейронів)? (б)~Стенд записує $1024$ електроди з частотою $20$~кГц по $10$~біт на відлік. У скільки разів передавання цього потоку при $1$~нДж на біт дорожче за потужність самої культури?
\end{exercise}

\begin{exercise}[\lvlB]
\label{ex:21:3}
\emph{Проєктування: придушити залпи.} Рідкісна стимуляція через багато електродів змушує кожен синапс викидати медіатор із частотою $r_b$; запас $\tau_{\text{відн}}\,\dd x/\dd t=1-x-Ur_b\tau_{\text{відн}}x$, $U=0.5$, $\tau_{\text{відн}}=0.8$~с. За якої найменшої $r_b$ запас не доростає до $x^*=0.9$; $0.99$ і наскільки при цьому слабшає синапс?
\end{exercise}

\begin{exercise}[\lvlB]
\label{ex:21:4}
\emph{Пам'ять резервуара не більша за число нейронів.} Резервуар з $N$ виходами $\mathbf r(t)$ отримує білий вхід $u(t)$ з нульовим середнім і одиничною дисперсією; $m_k$ --- найбільший квадрат кореляції зчитування $\mathbf w_k\cdot\mathbf r(t)$ з $u(t-k)$. Доведіть, що ємність пам'яті $\mathrm{MC}=\sum_{k\ge0}m_k\le N$.
\end{exercise}

\begin{exercise}[\lvlB]
\label{ex:21:5}
\emph{Моделювання: резервуар у кремнії.} Резервуар $\mathbf r(t+1)=\tanh(W\mathbf r(t)+\mathbf v\,u(t)+\mathbf c)$, $N=30$, $W$ випадкова зі спектральним радіусом $0.9$, $v_i\sim U(-0.5,0.5)$, $c_i\sim U(-0.2,0.2)$, $u\sim U(-1,1)$, зчитування --- гребенева регресія. Знайдіть ємність пам'яті задачі~\ref{ex:21:4} з $\tanh$ і без нього. Який резервуар пам'ятає більше?
\end{exercise}

\begin{exercise}[\lvlB]
\label{ex:21:6}
\emph{Зчитування для живого резервуара.} Органоїд дає $1024$ канали і $P=240$ навчальних записів двох класів. (а)~Скільки ознак $n$ можна залишити, щоб за формулою задачі~\ref{ex:19:3} випадкова розмітка була роздільною з імовірністю не більше $1\%$? (б)~На скільки сигм вища за вгадування точність $78\%$ на $100$ перевірних записах?
\end{exercise}

\begin{exercise}[\lvlC]
\label{ex:21:7}
\emph{Дослідження: учитель без широкомовного каналу.} Запропонуйте протокол стимуляції, що реалізує на стенді правило трьох множників із розділу~\ref{sec:dofa} через $8$--$1024$ електродів, і контроль із замороженим зчитуванням, який відрізняє навчання ваг від зсуву стану мережі. Який результат спростував би гіпотезу про навчання?
\end{exercise}
